% Generated draft: requires mathematical and source review.
\documentclass[11pt]{article}
\usepackage[margin=25mm]{geometry}
\usepackage{fontspec}
\setmainfont{Latin Modern Roman}
\usepackage{amsmath,amssymb,mathtools,unicode-math}
\setmathfont{Latin Modern Math}
\usepackage{enumitem,needspace,xurl,hyperref}
\setlength{\parindent}{0pt}
\setlength{\parskip}{5pt}
\setlength{\emergencystretch}{3em}
\newcommand{\N}{\mathbb{N}}
\newcommand{\Z}{\mathbb{Z}}
\newcommand{\Q}{\mathbb{Q}}
\newcommand{\R}{\mathbb{R}}
\newcommand{\C}{\mathbb{C}}
\newcommand{\F}{\mathbb{F}}
\newcommand{\A}{\mathbb{A}}
\newcommand{\PP}{\mathbb{P}}
\newcommand{\E}{\mathbb{E}}
\newcommand{\eps}{\varepsilon}
\newcommand{\dd}{\,\mathrm{d}}
\begin{document}
\section*{UnsolvedMath ID: 3100010}
\textbf{Source problem code:} catalogue entry
\subsection{Definitions and mathematical statement}
A regular tetrahedron is the closed convex hull of four points in $\mathbb R^3$ at equal pairwise distances. Fix one such tetrahedron $T$; congruent copies are its rigid-motion images, with no scaling. A conventional nonoverlapping common-point formulation seeks copies $T_1,\ldots,T_N$ and $p\in\mathbb R^3$ such that
\[
p\in\bigcap_{i=1}^{N}\partial T_i,\qquad
\operatorname{int}(T_i)\cap\operatorname{int}(T_j)=\varnothing
\quad(i\ne j).
\]
Boundary contacts away from $p$ are allowed by this particular formulation.

\textbf{Question.} What is the largest number of congruent regular tetrahedra that can touch at one common point? The source reports an arrangement with at least $20$ tetrahedra and an upper bound $22$, and asks for the exact value. Its research summary repeats
\[
20\leq N_{\max}\leq22
\]
and describes the upper bound as a simple volume bound, without providing a construction or proof.

The words touch at a point do not specify whether $p$ must be a vertex of every tetrahedron, whether contacts along additional edges or faces are allowed, or whether the pairwise intersections must equal $\{p\}$. Those stricter formulations are different from the displayed boundary-contact condition and are not silently imposed. The intended packing model must be matched to the reported bounds. There is no distinguished central tetrahedron: all copies are required to meet the same point, unlike a general central-body kissing-number problem. The supplied status remains open/triage rather than a verified exact value.
\subsection{Short English statement}
Find the maximum number of nonoverlapping congruent regular tetrahedra meeting at one point, between the reported bounds $20$ and $22$. The source needs an exact common-point and boundary-contact convention.
\subsection{Sources}
\begin{itemize}
\item Cooper, \emph{Combinatorial Problems I Like}, Geometry, source-order bullet 10, October 2020 snapshot, \url{https://people.math.sc.edu/cooper/combprob.html}; HTML extraction, accessed 2026-07-29. Later resolution was not checked item by item according to the record.
\item The supplied summary repeats the bounds $20$--$22$, mentions unspecified recent surveys, and contrasts a cube analogue with value $8$, but gives no separate construction, upper-bound, or survey citation.
\item Local registry \texttt{data/unsolvedmath/problems.json}, permanent ID $3100010$, source code catalogue entry, \emph{How many congruent regular tetrahedra can touch at a point}. The record labels this open and OPEN-TRIAGE, with NEEDS\_REVIEW source-status and rights qualifications.
\end{itemize}
\end{document}
